
\documentclass{article}
\usepackage{colortbl}
\usepackage{makecell}
\usepackage{multirow}
\usepackage{supertabular}

\begin{document}

\newcounter{utterance}

\centering \large Interaction Transcript for game `wordle', experiment `high\_frequency\_words\_no\_clue\_no\_critic', episode 10 with reflexion{-}Qwen3.5{-}9B.
\vspace{24pt}

{ \footnotesize  \setcounter{utterance}{1}
\setlength{\tabcolsep}{0pt}
\begin{supertabular}{c@{$\;$}|p{.15\linewidth}@{}p{.15\linewidth}p{.15\linewidth}p{.15\linewidth}p{.15\linewidth}p{.15\linewidth}}
    \# & \multicolumn{2}{c}{Player} && \multicolumn{2}{c}{Game Master} \\
    \hline

    \theutterance \stepcounter{utterance}  
    & & & \multicolumn{4}{p{0.6\linewidth}}{
        \cellcolor[rgb]{0.9,0.9,0.9}{
            \makecell[{{p{\linewidth}}}]{
                \texttt{\tiny{[P1$\langle$GM]}}
                \texttt{You have attempted this same task before and did not succeed. These are the plans you wrote after those attempts.} \\
\texttt{Trial \#1: New plan: The previous attempt was aborted because my explanation became excessively long and rambling, violating the instruction to provide a "concise" explanation. Additionally, I failed to finalize a valid guess within the explanation before the text cut off. Next time, I will strictly limit the explanation to one or two sentences that directly map the feedback to the specific letter positions of the new word, ensuring the resulting word is a valid five{-}letter English word that adheres to all constraints (e.g., 'c' at start, 'e' at position 4, 'r' not at position 2) before submitting the "guess:" tag.} \\
\\ 
\\ 
\texttt{You are a language wizard who likes to guess words by using the given rules.} \\
\\ 
\texttt{Welcome to Wordle! You have six attempts to guess the target word, a valid English word of five lowercase letters (a{-}z). Please use the tags "explanation:" and "guess:" to provide a concise explanation for each guess.} \\
\\ 
\texttt{For instance, if your guess is "apple", your response should be} \\
\texttt{explanation: this is a common five{-}letter English word, and I am starting my guess with this word.} \\
\texttt{guess: apple} \\
\\ 
\texttt{After each guess, your answer will be validated, and you will receive feedback indicating which letters are correct (green), which letters are correct but in the wrong position (yellow), and which letters are incorrect (red). This feedback can be useful in determining which letters to include or exclude in your next guess.} \\
\\ 
\texttt{For example, the feedback for "apple" might be:} \\
\texttt{guess\_feedback: a<yellow> p<yellow> p<green> l<yellow> e<red>} \\
\\ 
\texttt{The explanation should contain details about how the guess\_feedback is used to arrive at a new guess.} \\
\\ 
\texttt{Do not generate any other text.} \\
\\ 
\texttt{Let's begin with your first guess.} \\
            }
        }
    }
     \\ \\

    \theutterance \stepcounter{utterance}  
    & \multicolumn{4}{p{0.6\linewidth}}{
        \cellcolor[rgb]{0.561,0.737,0.561}{
            \makecell[{{p{\linewidth}}}]{
                \texttt{\tiny{[P1$\rangle$GM]}}
                \texttt{explanation: I will start with a common five{-}letter word containing frequent letters to maximize information gain.} \\
\texttt{guess: slake} \\
            }
        }
    }
    & & \\ \\

    \theutterance \stepcounter{utterance}  
    & & & \multicolumn{2}{p{0.3\linewidth}}{
        \cellcolor[rgb]{0.9,0.9,0.9}{
            \makecell[{{p{\linewidth}}}]{
                \texttt{\tiny{[GM$|$GM]}}
                \texttt{attempts = 1} \\
\texttt{target\_word = cover} \\
\texttt{guess = slake} \\
\texttt{guess\_feedback = s<red> l<red> a<red> k<red> e<yellow>} \\
            }
        }
    }
    & & \\ \\

    \theutterance \stepcounter{utterance}  
    & & & \multicolumn{4}{p{0.6\linewidth}}{
        \cellcolor[rgb]{0.9,0.9,0.9}{
            \makecell[{{p{\linewidth}}}]{
                \texttt{\tiny{[P1$\langle$GM]}}
                \texttt{guess\_feedback: s<red> l<red> a<red> k<red> e<yellow>} \\
\\ 
\texttt{Provide your response only in this format.} \\
\texttt{explanation: details} \\
\texttt{guess: word} \\
            }
        }
    }
     \\ \\

    \theutterance \stepcounter{utterance}  
    & \multicolumn{4}{p{0.6\linewidth}}{
        \cellcolor[rgb]{0.561,0.737,0.561}{
            \makecell[{{p{\linewidth}}}]{
                \texttt{\tiny{[P1$\rangle$GM]}}
                \texttt{explanation: Since 'e' is yellow (position 5) and all other letters are red, the target word ends in 'e' and does not contain s, l, a, or k.} \\
\texttt{guess: crane} \\
            }
        }
    }
    & & \\ \\

    \theutterance \stepcounter{utterance}  
    & & & \multicolumn{2}{p{0.3\linewidth}}{
        \cellcolor[rgb]{0.9,0.9,0.9}{
            \makecell[{{p{\linewidth}}}]{
                \texttt{\tiny{[GM$|$GM]}}
                \texttt{attempts = 2} \\
\texttt{target\_word = cover} \\
\texttt{guess = crane} \\
\texttt{guess\_feedback = c<green> r<yellow> a<red> n<red> e<yellow>} \\
            }
        }
    }
    & & \\ \\

    \theutterance \stepcounter{utterance}  
    & & & \multicolumn{4}{p{0.6\linewidth}}{
        \cellcolor[rgb]{0.9,0.9,0.9}{
            \makecell[{{p{\linewidth}}}]{
                \texttt{\tiny{[P1$\langle$GM]}}
                \texttt{guess\_feedback: c<green> r<yellow> a<red> n<red> e<yellow>} \\
\\ 
\texttt{Provide your response only in this format.} \\
\texttt{explanation: details} \\
\texttt{guess: word} \\
            }
        }
    }
     \\ \\

    \theutterance \stepcounter{utterance}  
    & \multicolumn{4}{p{0.6\linewidth}}{
        \cellcolor[rgb]{0.561,0.737,0.561}{
            \makecell[{{p{\linewidth}}}]{
                \texttt{\tiny{[P1$\rangle$GM]}}
                \texttt{explanation: 'c' is correct at the start, 'r' and 'e' are correct but misplaced (so 'r' is not 2nd and 'e' is not 5th), and 'a', 'n' are incorrect. Since 'e' is yellow, it must be in position 2, 3, or 4. Given 'r' is yellow (not 2nd), and we need a word starting with 'c' ending with a vowel other than 'e' or having 'e' in the middle, let's try "creed" but 'd' is unknown. Wait, 'e' is yellow, so it's not at the end. Let's try "crept" where 'e' is at 2? No, 'r' is yellow meaning it's not at 2. If 'e' is yellow, it's not at 5. Let's try "crane" again? No, 'a' and 'n' were red. How about "creed"? 'e' at 2 and 5? No, only one 'e' in feedback? The feedback shows e<yellow>, meaning the 'e' in crane is in the wrong spot. So 'e' is not at 5. It could be at 2, 3, or 4. 'r' is yellow, so not at 2. 'c' is green. Let's try "crept": c{-}green, r{-}yellow (ok if r is at 2? No, r is yellow in crane, meaning r is not at 2. In crept, r is at 2. That would make r green if it was the right spot. But we don't know where r is. Wait, in crane, r was at 2 and got yellow, so r is not at 2. In crept, r is at 2, so that would be wrong. Let's try "credo": c{-}green, r{-}yellow (if r is not 2), e{-}yellow (if e is not 5), d{-}red?, o{-}red?. We don't know d or o. Let's try a word where e is at 3 or 4. "creed" has e at 2 and 5. If the target has only one e, then creed might be wrong. Let's try "crisp"? No, p is unknown. Let's try "crumb"? No. How about "crept"? c{-}green, r{-}yellow (so r is not 2), e{-}yellow (so e is not 5), p{-}red?, t{-}red?. But we don't know p or t. Let's try "crane" was bad. Let's try "credo". Or "creed". Wait, the target has c at 1, r somewhere, e somewhere (not 5), and} \\
            }
        }
    }
    & & \\ \\

    \theutterance \stepcounter{utterance}  
    & & & \multicolumn{2}{p{0.3\linewidth}}{
        \cellcolor[rgb]{0.9,0.9,0.9}{
            \makecell[{{p{\linewidth}}}]{
                \texttt{\tiny{[GM$|$GM]}}
                \texttt{Error: The response should contain the 'guess:' keyword exactly once.} \\
            }
        }
    }
    & & \\ \\

    \theutterance \stepcounter{utterance}  
    & & & \multicolumn{2}{p{0.3\linewidth}}{
        \cellcolor[rgb]{0.9,0.9,0.9}{
            \makecell[{{p{\linewidth}}}]{
                \texttt{\tiny{[GM$|$GM]}}
                \texttt{game\_result = ABORT} \\
            }
        }
    }
    & & \\ \\

    \theutterance \stepcounter{utterance}  
    & & \multicolumn{4}{p{0.6\linewidth}}{
        \cellcolor[rgb]{0.95,0.95,0.95}{
            \makecell[{{p{\linewidth}}}]{
                \texttt{\tiny{[P1$\leftrightarrow$P]}}
                \texttt{Post{-}episode Reflexion (not sent to GM):} \\
\texttt{New plan: The previous attempt was aborted because the explanation became excessively long, rambling, and failed to conclude with a valid guess before the text cut off, violating the "concise" requirement and the format rule. Next time, I will strictly limit the explanation to one or two sentences that directly map the feedback (c=green, r/e=yellow, a/n=red) to the specific letter positions of the new word, ensuring the resulting word is a valid five{-}letter English word that adheres to all constraints (e.g., 'c' at start, 'e' at position 4, 'r' not at position 2) before immediately submitting the "guess:" tag.} \\
            }
        }
    }
    & \\ \\

\end{supertabular}
}

\end{document}
